\documentclass[10pt,a4paper]{amsart}
\usepackage{amsmath,amssymb,amsthm,graphicx}
\usepackage{verbatim}
\newenvironment{reference}{\comment}{\endcomment}
%\newenvironment{reference}{}{}
\newenvironment{slogan}{\comment}{\endcomment}
\newenvironment{history}{\comment}{\endcomment}

% For commutative diagrams we use Xy-pic
\usepackage[all]{xy}

% We use 2cell for 2-commutative diagrams.
\xyoption{2cell}
% Theorem environments.
%
\theoremstyle{plain}
\newtheorem{theorem}[subsection]{Theorem}
\newtheorem{proposition}[subsection]{Proposition}
\newtheorem{lemma}[subsection]{Lemma}

\theoremstyle{definition}
\newtheorem{definition}[subsection]{Definition}
\newtheorem{example}[subsection]{Example}
\newtheorem{exercise}[subsection]{Exercise}
\newtheorem{situation}[subsection]{Situation}

\theoremstyle{remark}
\newtheorem{remark}[subsection]{Remark}
\newtheorem{remarks}[subsection]{Remarks}

\numberwithin{equation}{subsection}

% Macros
%
\def\lim{\mathop{\mathrm{lim}}\nolimits}
\def\colim{\mathop{\mathrm{colim}}\nolimits}
\def\Spec{\mathop{\mathrm{Spec}}}
\def\Hom{\mathop{\mathrm{Hom}}\nolimits}
\def\Ext{\mathop{\mathrm{Ext}}\nolimits}
\def\SheafHom{\mathop{\mathcal{H}\!\mathit{om}}\nolimits}
\def\SheafExt{\mathop{\mathcal{E}\!\mathit{xt}}\nolimits}
\def\Sch{\mathit{Sch}}
\def\Mor{\mathop{\mathrm{Mor}}\nolimits}
\def\Ob{\mathop{\mathrm{Ob}}\nolimits}
\def\Sh{\mathop{\mathit{Sh}}\nolimits}
\def\NL{\mathop{N\!L}\nolimits}
\def\CH{\mathop{\mathrm{CH}}\nolimits}
\def\proetale{{pro\text{-}\acute{e}tale}}
\def\etale{{\acute{e}tale}}
\def\QCoh{\mathit{QCoh}}
\def\Ker{\mathop{\mathrm{Ker}}}
\def\Im{\mathop{\mathrm{Im}}}
\def\Coker{\mathop{\mathrm{Coker}}}
\def\Coim{\mathop{\mathrm{Coim}}}

% Boxtimes
%
\DeclareMathSymbol{\boxtimes}{\mathbin}{AMSa}{"02}

%
% Macros for moduli stacks/spaces
%
\def\QCohstack{\mathcal{QC}\!\mathit{oh}}
\def\Cohstack{\mathcal{C}\!\mathit{oh}}
\def\Spacesstack{\mathcal{S}\!\mathit{paces}}
\def\Quotfunctor{\mathrm{Quot}}
\def\Hilbfunctor{\mathrm{Hilb}}
\def\Curvesstack{\mathcal{C}\!\mathit{urves}}
\def\Polarizedstack{\mathcal{P}\!\mathit{olarized}}
\def\Complexesstack{\mathcal{C}\!\mathit{omplexes}}
% \Pic is the operator that assigns to X its picard group, usage \Pic(X)
% \Picardstack_{X/B} denotes the Picard stack of X over B
% \Picardfunctor_{X/B} denotes the Picard functor of X over B
\def\Pic{\mathop{\mathrm{Pic}}\nolimits}
\def\Picardstack{\mathcal{P}\!\mathit{ic}}
\def\Picardfunctor{\mathrm{Pic}}
\def\Deformationcategory{\mathcal{D}\!\mathit{ef}}
\usepackage{cleveref}
\crefname{originalsection}{original section}{original sections}
\Crefname{originalsection}{Original section}{Original sections}
\makeatletter
\newlabel{duality-example-iso-on-RSheafHom-noetherian}{{0GEW}{source}{}{sourcepointer-1}{}}
\newlabel{duality-example-iso-on-RSheafHom-noetherian@cref}{{[originalsection][0][]0GEW}{[1][0][]source}{}{}{}}
\makeatother
\crefname{originallemma}{original lemma}{original lemmas}
\Crefname{originallemma}{Original lemma}{Original lemmas}
\makeatletter
\newlabel{duality-lemma-duality-bootstrap}{{0AUE}{source}{}{sourcepointer-2}{}}
\newlabel{duality-lemma-duality-bootstrap@cref}{{[originallemma][0][]0AUE}{[1][0][]source}{}{}{}}
\makeatother
\crefname{originalsection}{original section}{original sections}
\Crefname{originalsection}{Original section}{Original sections}
\makeatletter
\newlabel{duality-section-duality}{{0AU3}{source}{}{sourcepointer-3}{}}
\newlabel{duality-section-duality@cref}{{[originalsection][0][]0AU3}{[1][0][]source}{}{}{}}
\makeatother
\crefname{originalsection}{original section}{original sections}
\Crefname{originalsection}{Original section}{Original sections}
\makeatletter
\newlabel{duality-section-trace}{{0AWG}{source}{}{sourcepointer-4}{}}
\newlabel{duality-section-trace@cref}{{[originalsection][0][]0AWG}{[1][0][]source}{}{}{}}
\makeatother
\crefname{originallemma}{original lemma}{original lemmas}
\Crefname{originallemma}{Original lemma}{Original lemmas}
\makeatletter
\newlabel{dualizing-lemma-special-case-local-duality}{{0AAK}{source}{}{sourcepointer-5}{}}
\newlabel{dualizing-lemma-special-case-local-duality@cref}{{[originallemma][0][]0AAK}{[1][0][]source}{}{}{}}
\makeatother
\crefname{originalproposition}{original proposition}{original propositions}
\Crefname{originalproposition}{Original proposition}{Original propositions}
\makeatletter
\newlabel{dualizing-proposition-matlis}{{08Z9}{source}{}{sourcepointer-6}{}}
\newlabel{dualizing-proposition-matlis@cref}{{[originalproposition][0][]08Z9}{[1][0][]source}{}{}{}}
\makeatother
\crefname{originalequation}{original equation}{original equations}
\Crefname{originalequation}{Original equation}{Original equations}
\makeatletter
\newlabel{equation-a-torsion}{{0AXK}{source}{}{sourcepointer-7}{}}
\newlabel{equation-a-torsion@cref}{{[originalequation][0][]0AXK}{[1][0][]source}{}{}{}}
\makeatother
\crefname{originalreference}{original reference}{original references}
\Crefname{originalreference}{Original reference}{Original references}
\makeatletter
\newlabel{lemma-bound-a-torsion}{{0AXJ}{source}{}{sourcepointer-8}{}}
\newlabel{lemma-bound-a-torsion@cref}{{[originalreference][0][]0AXJ}{[1][0][]source}{}{}{}}
\makeatother
\crefname{originalreference}{original reference}{original references}
\Crefname{originalreference}{Original reference}{Original references}
\makeatletter
\newlabel{proposition-Grauert-Riemenschneider}{{0AXD}{source}{}{sourcepointer-9}{}}
\newlabel{proposition-Grauert-Riemenschneider@cref}{{[originalreference][0][]0AXD}{[1][0][]source}{}{}{}}
\makeatother
\title{A bound on the cokernel of the trace map for all modifications implies reduction to rational singularities}
\author{The Stacks Project Authors}
\date{stacks-project:resolve.tex:local-snapshot-20261012; GNU FDL 1.2 or later}
\begin{document}
\maketitle
\noindent\textit{Editorial note.} All delivered mathematical text is the author's own, copied from the pinned Stacks Project chapter file resolve.tex; no statement or proof step was rewritten, reordered or removed.\par\medskip
\noindent\textit{Editorial note.} The counted proof is the author's complete argument: the use of Grauert-Riemenschneider to make the trace map a map of modules on the base, the duality computation expressing its cokernel as a second Ext group of the first direct image, the identification of that group with a Matlis dual transporting the annihilation, and the final appeal to the bound on the torsion of the first cohomology.\par\medskip
\noindent\textit{Editorial note.} The label lemma-bound-dualizing-implies-bound names the counted statement, and the proof environment immediately following it in the source is delivered verbatim as its proof.\par\medskip
\noindent\textit{Editorial note.} The counted statement is carried by the author's own lemma environment in the chapter source.\par\medskip
\noindent\textit{Editorial note.} Pre-open gate: python operations/agent-authoring/tools/precard\_gate.py --source resolve.tex --label lemma-bound-dualizing-implies-bound --start 2232 --end 2288 --local-id 28087 returned RESULT=OK (L28087 free, occupancy FREE, public interval clear). Pre-assembly tag self-check: canonical tag resolved and not present in reference/public-items.jsonl.\par\medskip
\noindent\textit{Editorial note.} The statement and proof are copied from the pinned chapter file; the cited inputs are kept as external pointers with their canonical Stacks tags.\par\medskip
\noindent\textit{Allowed standard prerequisites (editorial).}\begin{itemize}
\item Grauert-Riemenschneider vanishing for the pushforward of the dualizing module. Same chapter, the proposition stating that the higher direct images of the dualizing module vanish, so that the trace map is a map on the zeroth pushforward. Supplies the identification of the derived pushforward of the dualizing module with its zeroth pushforward, making the trace map a map of modules on the base.
\item Duality for a modification and the distinguished triangle from the structure sheaf. Duality for Schemes, the section on trace maps and the lemmas identifying the dualizing complex of the modification and expressing the pushforward of the structure sheaf through duality. Supply the identity expressing the pushforward of the dualizing module as the derived hom of the pushforward of the structure sheaf and the distinguished triangle turned into the displayed short exact sequence.
\item Local duality and the Matlis equivalence on a complete local ring. Dualizing Complexes, the local duality lemma identifying the relevant Ext group with the Matlis dual and the proposition that Matlis duality is an equivalence of categories. Supplies the conclusion that the first direct image is annihilated by the same element as the Ext group, hence by the given element.
\item Identification of the first cohomology with the torsion of the cokernel. Same chapter, the lemma bounding the torsion of the first cohomology and the identity expressing that torsion as a cokernel. Supplies the final boundedness statement needed for reduction to rational singularities.
\end{itemize}\medskip
\ifdefined\hypertarget\hypertarget{sourcepointer-1}{}\fi
\noindent\textit{Source reference.} Original section 0GEW: sources/stacks-chow/tags.txt, source lines 20491--20491. Prerequisite pointer invoked at proof line(s) 2259; it names a standard fact of a sibling chapter that the counted argument uses, and every step of the counted proof built on it is retained verbatim in the delivered body.\par\medskip
\ifdefined\hypertarget\hypertarget{sourcepointer-2}{}\fi
\noindent\textit{Source reference.} Original lemma 0AUE: sources/stacks-chow/tags.txt, source lines 13647--13647. Prerequisite pointer invoked at proof line(s) 2257; it names a standard fact of a sibling chapter that the counted argument uses, and every step of the counted proof built on it is retained verbatim in the delivered body.\par\medskip
\ifdefined\hypertarget\hypertarget{sourcepointer-3}{}\fi
\noindent\textit{Source reference.} Original section 0AU3: sources/stacks-chow/tags.txt, source lines 13636--13636. Prerequisite pointer invoked at proof line(s) 2258; it names a standard fact of a sibling chapter that the counted argument uses, and every step of the counted proof built on it is retained verbatim in the delivered body.\par\medskip
\ifdefined\hypertarget\hypertarget{sourcepointer-4}{}\fi
\noindent\textit{Source reference.} Original section 0AWG: sources/stacks-chow/tags.txt, source lines 13719--13719. Prerequisite pointer invoked at proof line(s) 2251; it names a standard fact of a sibling chapter that the counted argument uses, and every step of the counted proof built on it is retained verbatim in the delivered body.\par\medskip
\ifdefined\hypertarget\hypertarget{sourcepointer-5}{}\fi
\noindent\textit{Source reference.} Original lemma 0AAK: sources/stacks-chow/tags.txt, source lines 12988--12988. Prerequisite pointer invoked at proof line(s) 2280; it names a standard fact of a sibling chapter that the counted argument uses, and every step of the counted proof built on it is retained verbatim in the delivered body.\par\medskip
\ifdefined\hypertarget\hypertarget{sourcepointer-6}{}\fi
\noindent\textit{Source reference.} Original proposition 08Z9: sources/stacks-chow/tags.txt, source lines 11367--11367. Prerequisite pointer invoked at proof line(s) 2282; it names a standard fact of a sibling chapter that the counted argument uses, and every step of the counted proof built on it is retained verbatim in the delivered body.\par\medskip
\ifdefined\hypertarget\hypertarget{sourcepointer-7}{}\fi
\noindent\textit{Source reference.} Original equation 0AXK: sources/stacks-resolve/pinned-original-source.tex, source lines 1965--1965. Prerequisite pointer invoked at proof line(s) 2287; it names a standard fact of the same chapter that the counted argument uses, and every step of the counted proof built on it is retained verbatim in the delivered body.\par\medskip
\ifdefined\hypertarget\hypertarget{sourcepointer-8}{}\fi
\noindent\textit{Source reference.} Original reference 0AXJ: sources/stacks-resolve/pinned-original-source.tex, source lines 1947--1947. Prerequisite pointer invoked at proof line(s) 2286; it names a standard fact of the same chapter that the counted argument uses, and every step of the counted proof built on it is retained verbatim in the delivered body.\par\medskip
\ifdefined\hypertarget\hypertarget{sourcepointer-9}{}\fi
\noindent\textit{Source reference.} Original reference 0AXD: sources/stacks-resolve/pinned-original-source.tex, source lines 1838--1838. Prerequisite pointer invoked at proof line(s) 2249; it names a standard fact of the same chapter that the counted argument uses, and every step of the counted proof built on it is retained verbatim in the delivered body.\par\medskip
\begin{lemma}
\label{lemma-bound-dualizing-implies-bound}
Let $A$ be a local normal Nagata domain of dimension $2$ which has a
dualizing complex $\omega_A^\bullet$. If there exists a nonzero $d \in A$
such that for all normal modifications $X \to \Spec(A)$ the cokernel of the
trace map
$$
\Gamma(X, \omega_X) \to \omega_A
$$
is annihilated by $d$, then reduction to rational singularities
is possible for $A$.
\end{lemma}

\begin{proof}
For $X \to \Spec(A)$ as in the statement we have to bound
$H^1(X, \mathcal{O}_X)$. Let $\omega_X$ be the dualizing module
of $X$ as in the statement of Grauert-Riemenschneider
(Proposition \ref{proposition-Grauert-Riemenschneider}).
The trace map is the map $Rf_*\omega_X \to \omega_A$ described
in Duality for Schemes, Section \ref{duality-section-trace}.
By Grauert-Riemenschneider we have $Rf_*\omega_X = f_*\omega_X$
thus the trace map indeed produces a map $\Gamma(X, \omega_X) \to \omega_A$.
By duality we have $Rf_*\omega_X = R\Hom_A(Rf_*\mathcal{O}_X, \omega_A)$
(this uses that $\omega_X[2]$ is the dualizing complex on $X$
normalized relative to $\omega_A[2]$,
see Duality for Schemes, Lemma \ref{duality-lemma-duality-bootstrap}
or more directly Section \ref{duality-section-duality} or even more directly
Example \ref{duality-example-iso-on-RSheafHom-noetherian}).
The distinguished triangle
$$
A \to Rf_*\mathcal{O}_X \to R^1f_*\mathcal{O}_X[-1] \to A[1]
$$
is transformed by $R\Hom_A(-, \omega_A)$ into the short exact sequence
$$
0 \to f_*\omega_X \to \omega_A \to
\Ext_A^2(R^1f_*\mathcal{O}_X, \omega_A) \to 0
$$
(and $\Ext_A^i(R^1f_*\mathcal{O}_X, \omega_A) = 0$ for $i \not = 2$;
this will follow from the discussion below as well).
Since $R^1f_*\mathcal{O}_X$ is supported in $\{\mathfrak m\}$, the
local duality theorem tells us that
$$
\Ext_A^2(R^1f_*\mathcal{O}_X, \omega_A) =
\Ext_A^0(R^1f_*\mathcal{O}_X, \omega_A[2]) =
\Hom_A(R^1f_*\mathcal{O}_X, E)
$$
is the Matlis dual of $R^1f_*\mathcal{O}_X$ (and the other
ext groups are zero), see
Dualizing Complexes, Lemma \ref{dualizing-lemma-special-case-local-duality}.
By the equivalence of categories inherent in Matlis duality
(Dualizing Complexes, Proposition \ref{dualizing-proposition-matlis}),
if $R^1f_*\mathcal{O}_X$ is not annihilated by $d$,
then neither is the $\Ext^2$ above. Hence we see that
$H^1(X, \mathcal{O}_X)$ is annihilated by $d$. Thus the required
boundedness follows from Lemma \ref{lemma-bound-a-torsion} and
(\ref{equation-a-torsion}).
\end{proof}

\end{document}
